\documentclass[12pt]{article}

\usepackage{sbc-template}
\usepackage{graphicx,url}
\usepackage[utf8]{inputenc}
\usepackage[brazil]{babel}
\usepackage{booktabs}

\sloppy

\title{DeepForensics: detecção paralela e distribuída de adulteração de
  imagens no CASIA 2.0 por múltiplas análises forenses e Random Forest}

\author{Caio Gomes Alcântara Gloria\inst{1}, Daniel Salgado Magalhães\inst{1},
  Felipe Silva Faria\inst{1},\\ João Lucas de Melo Quintão\inst{1},
  João Paulo de Castro\inst{1}, Lucas Araujo Barduino Rodrigues\inst{1}}

\address{Curso de Ciência da Computação\\
  Pontifícia Universidade Católica de Minas Gerais (PUC Minas)\\
  Belo Horizonte -- MG -- Brasil
}

\begin{document}

\maketitle

\begin{abstract}
  Editing tools have made manipulated images common, which motivates automatic
  detection methods. This paper describes DeepForensics, a system that labels
  images from the CASIA 2.0 dataset as authentic or tampered. Detection combines
  four classical passive cues, each turned into a calibrated 0--100 score:
  Error Level Analysis (ELA), noise inconsistency, perspective consistency via
  vanishing points, and relative object scale. A Random Forest learns to fuse
  the four scores into the final decision, replacing a hand-tuned weighted sum.
  The per-image workload is parallelized over the cores of the machine with the
  Python \texttt{multiprocessing} library, and we measure the gain through strong
  and weak scaling experiments. The system is also designed to run distributed
  across machines linked by a WireGuard overlay under a chief/worker model over
  XML-RPC, in which only image paths and scores would travel the network; we
  report the implemented single-machine parallel system and describe the
  distributed layer as planned work. A web interface (Flask plus HTML/JS) drives
  single-image analysis with visualizations, labelled-dataset evaluation, Random
  Forest training, and the scalability benchmarks.
\end{abstract}

\begin{resumo}
  A facilidade de editar imagens com as ferramentas atuais tornou comum a
  circulação de conteúdo adulterado, o que motiva métodos automáticos de
  detecção. Este artigo descreve o DeepForensics, um sistema que classifica
  imagens do conjunto CASIA 2.0 como autênticas ou adulteradas. A detecção
  combina quatro vestígios passivos clássicos, cada um convertido em um score
  calibrado de 0 a 100: a análise de nível de erro (ELA), a inconsistência de
  ruído, a consistência de perspectiva por pontos de fuga e a escala relativa
  dos objetos. Um \textit{Random Forest} aprende a fundir os quatro scores na
  decisão final, no lugar de uma soma de pesos ajustada à mão. A carga por
  imagem é paralelizada entre os núcleos da máquina com a biblioteca
  \texttt{multiprocessing} do Python, e medimos o ganho por meio de experimentos
  de escalabilidade forte e fraca. O sistema também está projetado para executar
  de forma distribuída entre máquinas ligadas por uma rede WireGuard, num modelo
  \textit{chief}/\textit{worker} sobre XML-RPC, em que apenas caminhos de imagem
  e scores trafegariam pela rede; relatamos o que já roda (o paralelismo em uma
  máquina) e descrevemos a camada distribuída como trabalho planejado. Uma
  interface web (Flask com HTML/JS) opera a análise de imagem única com
  visualizações, a avaliação de conjuntos rotulados, o treino do \textit{Random
  Forest} e os \textit{benchmarks} de escalabilidade.
\end{resumo}

\section{Introdução}

Editar uma imagem hoje não exige conhecimento técnico. Aplicativos de celular
removem objetos, trocam rostos e ajustam cenas em poucos toques, e o resultado
costuma ser convincente o bastante para enganar quem olha sem desconfiar. Quando
uma imagem desse tipo aparece em uma notícia, num processo ou numa rede social,
saber se ela foi alterada deixa de ser curiosidade e passa a ter consequência.
Esse problema é bastante estudado na área forense de imagens, que procura
indícios de manipulação na própria imagem, sem depender de marca d'água ou de
metadados, que são fáceis de remover ou forjar \cite{farid2009,birajdar2013}.

O problema que tratamos é simples de evidenciar e trabalhoso de resolver: dada
uma imagem, decidir se ela é autêntica ou foi adulterada. Diferentes formas de
adulteração deixam vestígios diferentes, e nenhum sinal isolado é suficiente.
Quando uma parte de uma imagem é recortada, colada e salva de novo, ela passa a
ter um histórico de recompressão distinto do resto da cena, incoerência que a
análise de nível de erro (\textit{Error Level Analysis}, ELA) revela. A região
inserida também costuma trazer um padrão de ruído de sensor incompatível com a
vizinhança, quebrar a geometria de perspectiva da cena ou exibir um tamanho
relativo implausível para a posição que ocupa. Nosso detector lê justamente
esses quatro sinais --- ELA, ruído, pontos de fuga e escala --- e deixa que um
classificador \textit{Random Forest} aprenda a combiná-los na decisão final.

Há ainda uma questão de escala. O CASIA 2.0, que usamos como base, reúne cerca
de 12 mil imagens \cite{dong2013casia}. Rodar as quatro análises em todas elas,
em sequência e numa máquina só, é lento. Esse é exatamente o tipo de carga ---
uma mesma análise repetida sobre muitas imagens independentes --- que justifica
paralelizar o trabalho entre os núcleos de cada máquina e, num passo seguinte,
distribuí-lo entre várias máquinas, que é o foco de sistemas deste trabalho.

\subsection{Objetivos}

O objetivo geral é construir um sistema capaz de classificar as imagens do CASIA
2.0 como autênticas ou adulteradas, integrando processamento e análise de
imagens, computação paralela e computação distribuída.

Como objetivos específicos, definimos: (i) implementar a detecção de adulteração
por quatro análises passivas --- ELA, ruído, pontos de fuga e escala ---, cada
uma produzindo um score de 0 a 100; (ii) fundir esses scores com um classificador
\textit{Random Forest} treinado, no lugar de pesos fixos ajustados à mão; (iii)
paralelizar o processamento das imagens dentro de cada máquina e projetar a
execução distribuída entre várias máquinas conectadas por uma rede privada; (iv)
calibrar a saída em um score interpretável, em cinco faixas; e (v) medir tanto a
qualidade da classificação quanto o ganho de desempenho do paralelismo. Para
operar o sistema, construímos ainda uma interface web que reúne a análise de
imagens, o treino do modelo, a avaliação sobre conjuntos rotulados e a execução
dos \textit{benchmarks}.

\subsection{Justificativa}

Do ponto de vista prático, detectar adulteração é relevante para jornalismo,
perícia e moderação de conteúdo, e a literatura mostra que técnicas passivas
conseguem indícios úteis mesmo sem informação externa \cite{farid2009}. Do ponto
de vista do nosso curso, o trabalho serve como exercício de integração: em vez de
tratar processamento de imagens, sistemas distribuídos e paralelismo como
assuntos separados, ele força a pensar como as três coisas conversam num sistema
real. A própria característica da carga, uma mesma análise repetida sobre muitas
imagens independentes, é o caso de uso clássico que motiva uma arquitetura
paralela e distribuída, o que torna o problema um bom pretexto para aplicar esses
conceitos sem que eles pareçam desnecessários.

\section{Revisão da Literatura}

A detecção passiva de adulteração se apoia na ideia de que toda edição deixa
algum vestígio, mesmo que invisível a olho nu. Os levantamentos da área
organizam as técnicas por tipo de vestígio que exploram: artefatos de
compressão, inconsistências de ruído, traços de interpolação de cor, repetição
de regiões e inconsistências geométricas ou de iluminação
\cite{farid2009,birajdar2013}. Nosso detector toma emprestado quatro desses
sinais, descritos a seguir.

A análise de nível de erro (\textit{Error Level Analysis}, ELA), popularizada por
Krawetz \cite{krawetz2007}, recomprime a imagem em JPEG com qualidade conhecida e
observa a diferença em relação à original. Regiões editadas tendem a ter um nível
de erro distinto do restante, porque passaram por um histórico de compressão
diferente. Um sinal aparentado é o do \textit{JPEG ghost} \cite{farid2009}, que
revela inconsistências quando uma região foi salva com qualidade JPEG diferente
do entorno. O ELA é uma das quatro análises que combinamos.

A análise de ruído parte da premissa de que regiões de origens diferentes têm
variância de ruído incompatível com a vizinhança \cite{mahdian2009}; nós a
empregamos medindo o ruído local e a sua dispersão entre blocos da imagem. As
inconsistências geométricas e de iluminação \cite{farid2009} fundamentam nossa
terceira análise: numa cena com perspectiva, retas paralelas no mundo convergem
para um ponto de fuga, e desvios desse padrão sugerem inserção de conteúdo de
outra cena. O mesmo raciocínio de perspectiva sustenta a quarta análise, que
compara o tamanho relativo dos objetos com o esperado para a posição que ocupam.

Outras famílias clássicas atacam vestígios que não tratamos. A detecção de
copy-move cuida do caso em que um pedaço da imagem é copiado para outro lugar
dela mesma, e os métodos baseados em pontos-chave, como o SIFT \cite{lowe2004},
são robustos a rotação e a escala \cite{christlein2012}. Não implementamos
copy-move; ele aparece aqui apenas para situar nossas escolhas no panorama
forense.

Mais recentemente, métodos baseados em aprendizado profundo passaram a dominar os
\textit{benchmarks}. Bayar e Stamm \cite{bayar2016} propõem uma camada
convolucional que aprende a realçar resíduos de manipulação, Zhou et al.
\cite{zhou2018} combinam cor e ruído numa rede de duas vias, e Nagm et al.
\cite{nagm2024} mostram que uma CNN compacta alimentada com a imagem ELA atinge
cerca de 94\% de acurácia no CASIA 2.0. Optamos por um caminho diferente: em vez
de treinar uma rede profunda, que exige mais dados e GPU, extraímos quatro
descritores interpretáveis e os fundimos com um \textit{Random Forest}, um
ensemble de árvores de decisão. A escolha mantém o custo por imagem baixo, sem
GPU, e concentra o esforço na parte de sistemas, que é o objetivo do trabalho.

\subsection{Posicionamento deste trabalho}

A maior parte da literatura busca um detector melhor e o avalia isoladamente.
Nosso trabalho tem outro foco. Não propomos um detector novo: reaproveitamos
quatro análises passivas clássicas e as combinamos com um classificador padrão.
O que tratamos é o lado de sistemas, pouco discutido nesses artigos: como
orquestrar a análise sobre um conjunto grande de imagens, paralelizando-a entre
os núcleos e, em projeto, distribuindo-a entre máquinas, calibrando a saída e
medindo o ganho de desempenho. A escolha por descritores baratos é deliberada:
todas as quatro análises rodam em CPU, sem treino pesado, o que mantém o foco na
parte paralela e distribuída. A Tabela~\ref{tab:comparativo} resume esse
posicionamento.

\begin{table}[ht]
\centering
\caption{Comparação entre as abordagens da literatura e a deste trabalho}
\label{tab:comparativo}
\begin{tabular}{lccc}
\toprule
\textbf{Aspecto} & \textbf{ELA clássico} & \textbf{Rede profunda} & \textbf{Este trabalho} \\
\midrule
Sinais usados          & 1 (ELA)     & aprendidos   & 4 (ELA, ruído, VP, escala) \\
Treinamento de modelo  & não         & sim (GPU)    & sim (Random Forest, leve) \\
Interpretabilidade     & alta        & baixa        & média/alta \\
Foco em sistemas       & raro        & raro         & central \\
Custo por imagem       & baixo       & alto (GPU)   & baixo (CPU) \\
\bottomrule
\end{tabular}
\end{table}

\section{Metodologia}

\subsection{Conjunto de dados}

Usamos o CASIA 2.0 \cite{dong2013casia}, que reúne 7.491 imagens autênticas (Au)
e 5.123 imagens adulteradas (Tp), nos formatos JPEG, BMP e TIFF. As adulterações
são de dois tipos: \textit{splicing}, que cola conteúdo de outra imagem, e
copy-move, que duplica regiões da mesma imagem. O conjunto traz o rótulo de cada
imagem, o que permite avaliação supervisionada. O repositório do projeto embarca
uma amostra reduzida (50~Au e 50~Tp) para execuções rápidas, e o sistema aponta
automaticamente para as pastas \texttt{Dataset/Au} e \texttt{Dataset/Tp} com o
conjunto completo quando disponível. Para treinar o classificador, extraímos as
\textit{features} de todas as imagens e usamos uma partição estratificada em 75\%
treino e 25\% teste, complementada por validação cruzada de cinco dobras, sem
nunca treinar o modelo com imagens do conjunto de teste.

\subsection{Análises forenses}

Toda a detecção parte das mesmas quatro análises, calculadas de forma
independente sobre a imagem. Cada uma devolve um score de 0 a 100, em que valores
altos indicam consistência (autenticidade) e valores baixos, indícios de
manipulação.

\begin{itemize}
  \item \textbf{ELA (nível de erro).} A imagem é recomprimida em JPEG
    (qualidade base 95 e depois 75) e tomamos o valor absoluto da diferença
    pixel a pixel. Sobre esse mapa, em blocos de $32\times32$, medimos a média e
    o desvio do erro; média alta e dispersão alta entre blocos penalizam o score.
    Como a recompressão é interna, a análise se aplica também a BMP e TIFF.
  \item \textbf{Ruído.} Estimamos o ruído como a diferença entre a imagem e uma
    versão suavizada por filtro de mediana. Calculamos a variância do ruído em
    blocos de $32\times32$ e o coeficiente de variação (desvio sobre média) entre
    esses blocos; um coeficiente alto indica regiões com ruído incompatível e
    reduz o score \cite{mahdian2009}.
  \item \textbf{Pontos de fuga (VP).} Detectamos retas com Canny e a transformada
    de Hough probabilística e estimamos o ponto de fuga pela mediana das
    interseções de muitos pares de retas sorteados, num esquema do tipo RANSAC.
    Medimos então a divergência angular média das retas em relação a esse ponto;
    perspectiva quebrada baixa o score.
  \item \textbf{Escala/blob.} Após operações morfológicas, extraímos os maiores
    contornos (objetos) e comparamos a razão real entre as alturas com a razão
    esperada pela posição vertical de cada um na cena (perspectiva). Divergências
    grandes na razão logarítmica são marcadas como anomalias e reduzem o score.
\end{itemize}

\subsection{Fusão por Random Forest}

A primeira versão do detector combinava os quatro scores por pesos adaptativos ---
mais peso para ELA e ruído em retratos (poucas retas) e mais peso para pontos de
fuga em cenas arquitetônicas (muitas retas). Esses pesos eram fixados à mão.
Substituímos essa combinação por um classificador \textit{Random Forest} treinado,
que recebe como entrada os quatro scores e aprende a fronteira de decisão a partir
dos rótulos Au=0 (autêntica) e Tp=1 (adulterada). O modelo é treinado fora de
linha, persistido em disco em formato \texttt{.pkl}, e carregado na inferência.
Para uma imagem nova enviada por \textit{upload}, as quatro análises rodam, o
\textit{Random Forest} devolve a probabilidade de adulteração, e o sistema reporta
a classe e a probabilidade. A fusão deixa de ser ajustada à mão e passa a ser
guiada por dados, fácil de retreinar quando o conjunto cresce; a soma de pesos
antiga permanece apenas como alternativa quando ainda não há modelo treinado.

\subsection{Calibração e decisão}

A saída do detector é um score de 0 a 100, mapeado em cinco faixas por limiares
fixos: até 20, MANIPULADA; até 40, ALTA CHANCE DE MANIPULAÇÃO; até 60,
INCONCLUSIVA; até 80, CONSISTÊNCIA MÉDIA; acima de 80, CONSISTENTE. Com o
\textit{Random Forest}, o score é derivado da probabilidade de autenticidade
($(1-P_{\text{adulterada}})\times 100$), de modo que a mesma escala de cinco
faixas continua válida, e a decisão binária autêntica/adulterada é a própria
predição do classificador. Reportamos precisão, revocação, F1 e acurácia sobre o
conjunto de teste.

\subsection{Arquitetura distribuída (projetada)}

A camada distribuída está projetada, mas ainda não implementada na versão atual;
descrevemos aqui o desenho previsto. As máquinas se comunicariam por uma rede
privada montada com WireGuard \cite{donenfeld2017wireguard}. Como as máquinas do
grupo ficam atrás de NAT, sem IP público, usaríamos uma instância EC2 na AWS, com
IP fixo, como concentrador (\textit{hub}) da rede: o ponto de endereço estável ao
qual todas se conectam, na sub-rede \texttt{10.0.0.0/24}, e por onde o tráfego
entre elas é roteado. O hub não executaria lógica da aplicação; serviria apenas
de rota.

Sobre essa rede rodaria o sistema no modelo \textit{chief}/\textit{worker}. O
\textit{chief} particionaria a lista de imagens em fatias, uma por nó, e
despacharia a cada \textit{worker} a sua fatia. Cada \textit{worker} exporia um
servidor XML-RPC com as operações de \textit{ping} e análise de lote; o
\textit{chief} chamaria esses servidores em paralelo, com uma \textit{thread} por
nó, e agregaria os resultados à medida que voltassem. Uma decisão importante do
desenho é não trafegar imagens pela rede: como todo o tráfego passaria pelo hub,
mandar os bytes das imagens transformaria a banda da instância em gargalo. Por
isso cada \textit{worker} guardaria uma cópia local do CASIA 2.0, e pela rede
circulariam apenas os caminhos relativos das imagens e os scores de volta, que são
pequenos.

\subsection{Paralelismo dentro do nó}

Dentro de cada nó, a análise das imagens é paralelizada com a biblioteca
\texttt{multiprocessing} do Python, dimensionando o \textit{pool} pelo número de
núcleos da máquina. Usamos processos, e não \textit{threads}, porque o trabalho é
limitado por CPU e o GIL do Python impediria paralelismo real com
\textit{threads}; o método de início é \texttt{spawn}, padrão no Windows. Cada
imagem é processada de forma independente por uma função que roda as quatro
análises e devolve apenas os scores (sem as imagens de visualização), o que torna
o custo de serialização entre processos pequeno. Esse mesmo \textit{pool} é
reutilizado tanto na extração de \textit{features} para o treino quanto na
avaliação de conjuntos e nos \textit{benchmarks}.

\subsection{Metodologia de avaliação}

Avaliamos duas coisas distintas, que não devem ser misturadas: a qualidade da
detecção e o desempenho.

Para a detecção, reportamos precisão, revocação, F1 e acurácia do \textit{Random
Forest} sobre o conjunto de teste, com validação cruzada de cinco dobras, e
inspecionamos a importância de cada uma das quatro análises segundo o modelo.

Para o desempenho, no eixo do paralelismo dentro de um nó, medimos escalabilidade
forte (mesma carga, variando o número de \textit{workers} do \textit{pool}: 1, 2,
4 e o total de núcleos) e escalabilidade fraca (carga por \textit{worker} fixa,
crescendo dados e \textit{workers} juntos), interpretadas à luz das leis de Amdahl
\cite{amdahl1967} e de Gustafson \cite{gustafson1988}. Além de tempo e
\textit{speedup} ($T_1/T_N$), calculamos a eficiência ($\text{speedup}/N$) e um
termo de \textit{overhead} ($\max(0,\, T_N \cdot N - T_1)$), que estima o custo de
\textit{fork} e de serialização. No eixo distribuído, já projetado, a avaliação
prevista compara a execução em uma única máquina com a execução entre vários nós
sobre a mesma carga; a diferença mediria o \textit{overhead} de rede e
coordenação. Os \textit{benchmarks} de paralelismo são disparados tanto pela linha
de comando quanto pela interface web, que gera um painel com as curvas de
\textit{speedup}, eficiência, tempo e \textit{overhead}.

\section{Resultados}

A infraestrutura de detecção e de paralelismo está implementada e em uso. Estão
prontos: as quatro análises forenses (ELA, ruído, pontos de fuga e escala), a
combinação por pesos adaptativos como alternativa, o classificador \textit{Random
Forest} com treino, persistência em \texttt{.pkl} e predição por \textit{upload},
a API web em Flask (análise de imagem única com visualizações, lote, conjunto
rotulado com métricas, treino e \textit{status} do modelo, e \textit{benchmark}),
os \textit{benchmarks} de escalabilidade forte e fraca, e a interface web
(HTML/JS/CSS) com telas de \textit{dashboard}, envio de imagem, avaliação de
conjunto e painel de \textit{benchmark}. A camada distribuída multi-máquina
(WireGuard, hub EC2 e \textit{chief}/\textit{worker} sobre XML-RPC) está projetada,
mas ainda não implementada nesta versão.

A Tabela~\ref{tab:importancias} traz um resultado concreto já obtido: a
importância que o \textit{Random Forest} atribuiu a cada análise, medida numa
execução sobre a amostra embarcada (50~Au e 50~Tp). Os pontos de fuga e o ELA
foram os sinais mais informativos, seguidos pelo ruído; a escala contribuiu menos.

\begin{table}[ht]
\centering
\caption{Importância das quatro análises segundo o Random Forest (amostra reduzida)}
\label{tab:importancias}
\begin{tabular}{lc}
\toprule
\textbf{Análise} & \textbf{Importância} \\
\midrule
Pontos de fuga (VP) & 34{,}5\% \\
ELA                 & 29{,}6\% \\
Ruído               & 25{,}0\% \\
Escala              & 10{,}9\% \\
\bottomrule
\end{tabular}
\end{table}

A Tabela~\ref{tab:deteccao} mostra como reportamos a qualidade da detecção. Na
amostra reduzida do repositório, com apenas quatro \textit{features} e poucas
imagens, o \textit{Random Forest} ficou próximo do acaso (acurácia em torno de
50\% em validação cruzada) --- esperado nesse regime, em que os scores saturam e
há poucos exemplos. A avaliação definitiva, sobre o CASIA 2.0 completo e,
possivelmente, com \textit{features} adicionais (estatísticas brutas de ELA e
ruído), está em curso e seus valores serão consolidados na execução final.

\begin{table}[ht]
\centering
\caption{Métricas de detecção (a consolidar sobre o CASIA 2.0 completo)}
\label{tab:deteccao}
\begin{tabular}{lcccc}
\toprule
\textbf{Configuração} & \textbf{Precisão} & \textbf{Revocação} & \textbf{F1} & \textbf{Acurácia} \\
\midrule
Pesos adaptativos (referência)      & --- & --- & --- & --- \\
Random Forest (amostra reduzida)    & --- & --- & --- & $\sim$50\% \\
Random Forest (CASIA 2.0 completo)  & --- & --- & --- & --- \\
\bottomrule
\end{tabular}
\end{table}

A Tabela~\ref{tab:escalaforte} apresenta o formato da escalabilidade forte no
eixo de \textit{workers}, fixando o conjunto de imagens. O \textit{speedup}
esperado cresce com o número de \textit{workers}, mas com retorno decrescente
conforme a fração sequencial passa a dominar, como prevê a lei de Amdahl
\cite{amdahl1967}.

\begin{table}[ht]
\centering
\caption{Escalabilidade forte no eixo de \textit{workers}, carga fixa (a preencher)}
\label{tab:escalaforte}
\begin{tabular}{lcccc}
\toprule
\textbf{Workers} & \textbf{Tempo (s)} & \textbf{Speedup} & \textbf{Eficiência} & \textbf{Overhead (s)} \\
\midrule
1 & --- & 1{,}00 & 1{,}00 & 0{,}00 \\
2 & --- & --- & --- & --- \\
4 & --- & --- & --- & --- \\
$N$ núcleos & --- & --- & --- & --- \\
\bottomrule
\end{tabular}
\end{table}

A Tabela~\ref{tab:overhead} antecipa o experimento do eixo distribuído, ainda
planejado: a diferença entre a execução numa única máquina e a execução entre
vários nós isolaria o \textit{overhead} de serialização, roteamento pelo hub e
coordenação por \textit{threads}.

\begin{table}[ht]
\centering
\caption{Custo da distribuição: máquina única \textit{vs.} distribuída (planejado)}
\label{tab:overhead}
\begin{tabular}{lc}
\toprule
\textbf{Cenário} & \textbf{Tempo total (s)} \\
\midrule
Uma máquina (todos os núcleos)          & --- \\
Distribuído (vários \textit{workers})   & --- \\
\bottomrule
\end{tabular}
\end{table}

% TODO: ao consolidar os dados, acrescentar o painel de CP gerado pela interface
% (speedup vs Amdahl, eficiencia forte/fraca, tempo e overhead) como figura.

\section{Conclusões}

Este trabalho descreveu o DeepForensics, um sistema para classificar imagens do
CASIA 2.0 como autênticas ou adulteradas combinando processamento de imagens,
computação paralela e computação distribuída. A contribuição principal não está
em um detector novo, e sim na forma de orquestrar quatro análises passivas
clássicas --- ELA, ruído, pontos de fuga e escala --- fundidas por um
\textit{Random Forest}, sobre um sistema paralelo e, em projeto, distribuído, com
uma saída interpretável em cinco faixas e uma interface web de operação.

Algumas decisões se mostraram acertadas ainda na fase de projeto. Padronizar
análises baratas, todas em CPU, manteve o custo por imagem baixo e fez o sistema
cobrir todo o conjunto, inclusive os formatos sem perda. Trocar a soma de pesos
ajustada à mão por um \textit{Random Forest} deixou a fusão dos sinais guiada por
dados e trivial de retreinar. Processar cada imagem de forma independente, com um
\textit{pool} de processos que troca apenas scores, viabilizou o paralelismo sem
um custo grande de serialização.

O trabalho tem limitações que reconhecemos. A avaliação consolidada ainda depende
de rodar o conjunto completo: na amostra reduzida, com apenas quatro
\textit{features}, o classificador fica próximo do acaso, o que indica a
necessidade de mais dados e, provavelmente, de descritores adicionais. A camada
distribuída está desenhada, mas não implementada; quando estiver, sua partição de
carga será estática, dividida por igual, e não reequilibra sozinha quando os nós
têm desempenho diferente. Por fim, a gerência das chaves do WireGuard, prevista a
partir de um único ponto, não seria aceitável num cenário real, em que cada
máquina geraria a própria chave privada.

Como trabalho futuro, vale concluir a avaliação sobre o CASIA 2.0 completo,
enriquecer o vetor de \textit{features} com estatísticas brutas de ELA e ruído,
implementar de fato a camada \textit{chief}/\textit{worker} sobre a rede WireGuard
e substituir a partição estática por uma fila dinâmica de trabalho.

\bibliographystyle{sbc}
\bibliography{sbc-template}

\end{document}
